\section{案例 2：下班后无法开始锻炼}

\subsection*{情境}
一名员工每天早晨都认同规律锻炼的重要性，但在耗竭的工作日结束后，很少真正走进健身房。

\subsection*{框架映射}
\begin{itemize}[nosep]
  \item \textbf{当前状态：} 工作结束时的耗竭感，加上朝家里漂移的趋势。
  \item \textbf{目标状态：} 抵达健身房，并完成一套预先准备好的训练。
  \item \textbf{主导变量：} 低 $\Will$、升高的 $\DV$，以及环境依赖地重新锁进休息机制。
  \item \textbf{错过的决策窗口：} 进入家门之前的通勤边界。
\end{itemize}

\subsection*{机制解释}
这个案例展示了时机与路径依赖如何相互叠加。如果通勤几乎总是稳定终止于“回家休息”这条序列，那么在有意识权衡开始之前，轨迹其实已经被部分选定了。用高级动力学语言说，回家路线像一条重新落回盆地的通道。用涌现语言说，反复出现的“工作到家到瘫倒”序列，会沉淀为一种有历史巩固效应的转变策略。

\subsection*{证据车道纪律}
当前的研究主干支持一种\textbf{中等}解释：稳定性架构与转变时机确实重要。它\emph{并不}支持一种强主张，即晚间锻炼失败已经被映射到一个可清晰经验识别的特定神经吸引子上。这里更严格、也更有生物学根据的教训是：状态持续性、稳定性以及转变成本，本身就是严肃的系统性质，所以干预应该优先瞄准这些性质。

\subsection*{操作性修复}
\begin{enumerate}[nosep]
  \item 把决策前移：上班前就把运动包收好。
  \item 选定一套默认训练，而不是保留一整菜单的可选项。
  \item 让通勤路线经过健身房，而不是直接回家。
  \item 把成功定义为“物理进入场地并完成一段预先设定的最小训练模块”，而不是理想化的完整训练。
\end{enumerate}

\subsection*{迁移教训}
这个案例说明一条通用规则：如果当前路线在思考开始之前就已经把旧状态重建起来了，那么应当重设计路线，而不是要求一个已经耗竭的自我在原地跨过全部障碍。